\scanpage{51}
$E,G$ \hfill \uline{\unsure{Exemples|Exemple} des morphismes et des foncteurs de localisation}

$B_G\qquad H\lhook\joinrel\xrightarrow[i]{}G\qquad X=G/H\in\mathrm{Ob}\,B_G$ muni de $\varepsilon\in\Gamma(X)$

\diagram{\begin{tikzcd}[ampersand replacement=\&] (B_G)_{/X}\arrow[r,"\simeq","{\varphi}"']\arrow[d,"{j_X}"']\& B_H\arrow[d,"{B_i}"]\\ B_G\&B_G\end{tikzcd}}

\diagram{\begin{tikzcd}[ampersand replacement=\&] \deleted{\unsure{X\to B_G}}\&\\ E\arrow[r]\&X\\F\arrow[u]\arrow[r,dashed]\& e_E\arrow[u,"{\varepsilon}"']\end{tikzcd}}

\[\begin{aligned}\varphi(E)&=E\times_X(e_E,\varepsilon)\\\psi(F)&=G\wedge_H F\end{aligned}\]

\[\begin{aligned}&\deleted{\unsure{B_i^*}(E)}.\\&\varphi\bigl(j^*(E)\bigr)=(E\times X)\times_X(e_E,\varepsilon)\simeq E\end{aligned}\]

oubli partiel des opérations de $G$ au profit de $H$

Donc $j_X$, via l'équivalence $\varphi$, s'interprète comme $B_i$. On a

\[\begin{aligned}j_{X!}&\quad\text{s'interprète comme}\quad F\longmapsto G\wedge_H F\\j_{X*}&\quad\text{s'interprète comme}\quad F\mapsto\underline{\mathrm{Hom}}_H(G,F)\end{aligned}\]

plus généralement si $Z\in\mathrm{Ob}\,B_G$, les foncteurs \deleted{\unsure{``oubli}}

\diagram{\begin{tikzcd}[ampersand replacement=\&,column sep=small] (B_G)_{/(Z,G)}\arrow[r,"{\text{oubli brut}}"]\arrow[d,"{\text{oubli partiel }G}"']\&B_{\unsure{G|e}}\arrow[d,"{\text{oubli partiel }G}"]\arrow[r]\&\deleted{\unsure{B}}_H\\ (B_H)_{/(Z,H)}\arrow[r,"{\text{oubli brut}}"']\&B_H\&\end{tikzcd}}

sont des foncteurs de localisation

Cas des groupes discrets $(E=(\mathrm{Ens}))$ : Si $G$ est le groupe fondamental \inserted{en $z\in Z$} d'un espace top. $Z$ connexe, \inserted{loc. \unsure{conn|connexe}} et \inserted{loc\textsuperscript{t}} simplement connexe, $B_G\simeq\text{revêtements de }Z$, \inserted{théorie de galois pour \unsure{$Z$} rev. de $Z$} et la donnée de $X$ s'identifie à la donnée d'un revêtement \inserted{connexe} de $Z$, pointé au dessus de $\unsure{e|z}$ par un $x\in X$.

Alors la \deleted{\unsure{corresp.}} équivalence $\varphi$ \deleted{est} s'interprète \inserted{(qui est \unsure{théor.} de Galois \deleted{pour $Z$} $\operatorname{Rev}Z$)} comme l'équivalence entre $\operatorname{Rev}(X)$ (pour $(B_G)_{/X}$) et la \deleted{théorie de} $B_H$, donc c'est la théorie de Galois pour $X$ et le pt base $x$ (compte tenu que $\pi_1(X,x)\simeq H\subset G$)

TSVP
% archive sheet number: 50
